% practice.tex -- Module 6 (Electromagnetism): Practice set.
% 24 new HSC-style questions: Section I 8 multiple choice, Section II 10 short answer,
% Section III 6 extended response. 82 marks. Ordered from easy to HSC-hard.
%   Q24(a) 2 -> 3 marks, so that each calculation step earns a mark; total 80 -> 82.)
% Chapter tags: ch1 Q1, Q8, Q9, Q18, Q24; ch2 Q2, Q4, Q10, Q12, Q22;
%   ch3 Q6, Q7, Q13, Q17, Q21; ch4 Q3, Q11, Q16, Q19; ch5 Q5, Q14, Q15, Q20, Q23.
% All questions are new (new contexts and numbers). None is copied from a booklet, the
% homework (HW Electromagnetism 1-5), the mock papers (B1, B2), the 2025 HSC paper, or a
% worked example or Quick check of Chapters 1-5 of this book.
% Directions were checked with vectors (F = q v x B, F = I l x B): +x right, +y up the
% page, +z out of the page.
% \Needspace before each multiple-choice question keeps its stem and options together, and
% before each question that opens with a diagram keeps its header with the diagram
% (conditional, layout-independent; no \clearpage or \pagebreak).
% Local helpers (prefix "prac"), defined below: \pracmarks, \pracblank, TikZ styles
% pracdot, praccross, pracWire, pracN, pracS. (New macros for the practice set only.)
% Figures and tables here use \caption* (no number): this chapter is unnumbered, so a
% numbered caption would continue the numbering of Chapter 5.
% \pracmarks{...}: marks at the right end of the last line of an item; if they do not
% fit, they move to the right end of a new line (TeXbook trick, ends the paragraph).
\providecommand\pracmarks[1]{{\unskip\nobreak\hfil\penalty50\hskip1em\hbox{}\nobreak
  \hfill\mbox{\textcolor{hscmuted}{\small(#1)}}\parfillskip=0pt\finalhyphendemerits=0\par}}
% \pracblank: an empty cell for the student to fill in (a short grey line).
\providecommand\pracblank{\textcolor{hscmuted}{\rule[-2pt]{2.6em}{0.4pt}}}
\tikzset{
  pracdot/.pic={\fill[hscfield!65] (0,0) circle (1.5pt);},
  praccross/.pic={\draw[hscfield!65, line width=0.7pt] (-0.075,-0.075) -- (0.075,0.075)
    (-0.075,0.075) -- (0.075,-0.075);},
  pracWire/.style={hscink, line width=1.4pt},
  pracN/.style={fill=hscforce!15, draw=hscink, line width=0.5pt},
  pracS/.style={fill=hscvel!15, draw=hscink, line width=0.5pt},
}

\unnumberedchapter[\faDumbbell]{Practice set}

These 24 questions are written in the style of the HSC exam. They test all five chapters of
this book. They start easy and get harder.
\begin{itemize}
  \item \textbf{Section I:} 8 multiple-choice questions, 1 mark each. Choose A, B, C or D.
  \item \textbf{Section II:} 10 short-answer questions, 3 or 4 marks each.
  \item \textbf{Section III:} 6 extended-response questions, 5 to 8 marks each.
\end{itemize}
The total is 82 marks. Allow about 2 hours and 30 minutes. (In the HSC exam you have about
1.8 minutes for each mark.)

\textbf{For every question:}
\begin{itemize}
  \item Use the formula sheet and these data-sheet values: size of the charge on an electron
        or a proton $e = \qty{1.602e-19}{C}$, $m_\mathrm{e} = \qty{9.109e-31}{kg}$,
        $m_\mathrm{p} = \qty{1.673e-27}{kg}$, $\mu_0 = 4\pi \times \qty{e-7}{N.A^{-2}}$,
        $g = \qty{9.8}{m.s^{-2}}$.
  \item Ignore gravity for electrons, protons and ions. Ignore air resistance.
  \item Show your working. Give each answer with a unit and sensible significant figures.
        For a force, a field or a current, also give its direction when the question asks.
\end{itemize}

The marks for each part of a question are shown on the right, like this:
\textcolor{hscmuted}{\small(2 marks)}. The tag next to each question number (for example
\textcolor{hscbrand}{Ch~2}) shows which chapter the question tests. If you get a question
wrong, revise that chapter. Full worked answers and marking criteria are at the back of the
book. Click a question number to jump to its answer.

% =====================================================================
\Needspace{14\baselineskip}
\section*{Section I: Multiple choice}

\Needspace{10\baselineskip}
\begin{practiceq}{1}{1}\label{pq:platepath}
An electron enters the uniform electric field between two parallel charged plates. When it
enters, it moves parallel to the plates. Which statement about its motion between the plates
is correct?
\choices[1]{Its velocity does not change, because the field is uniform.}
  {Its path is a parabola, like the path of a ball thrown horizontally.}
  {Its path is an arc of a circle, because the force is always at \ang{90} to its velocity.}
  {Its velocity along the plates increases, because the field does work on it.}
\end{practiceq}

\Needspace{12\baselineskip}
\begin{practiceq}{1}{2}\label{pq:wiredir}
\begin{sidebyside}[0.30]
  \centering
  \begin{tikzpicture}[x=1cm, y=1cm]
    \foreach \x in {-1.0,-0.5,0.5,1.0}
      \foreach \y in {-1.0,-0.5,0,0.5,1.0}
        \pic at (\x,\y) {pracdot};
    \draw[hscink, line width=1.5pt, midarrow] (0,1.35) -- (0,-1.35);
    \node[font=\small] at (-0.24,0.25) {$I$};
    \node[diagramlabel, text=hscfield, above] at (0,1.4) {$B$ out of the page};
  \end{tikzpicture}
\nextside
A long, straight wire lies in the plane of the page. It carries a current $I$ down the page.
A uniform magnetic field points out of the page (see the diagram). Which way is the magnetic
force on the wire?
\choices[2]{to the left}{to the right}{up the page}{into the page}
\end{sidebyside}
\end{practiceq}

\Needspace{8\baselineskip}
\begin{practiceq}{1}{4}\label{pq:current}
An ideal transformer has 120 turns on its primary coil and 2400 turns on its secondary coil.
The current in the secondary coil is \qty{0.25}{A}. What is the current in the primary
coil?
\choices{\qty{0.0125}{A}}{\qty{0.25}{A}}{\qty{5.0}{A}}{\qty{100}{A}}
\end{practiceq}

\Needspace{11\baselineskip}
\begin{practiceq}{1}{2}\label{pq:newton3}
Two long, straight, parallel wires, P and Q, carry currents in opposite directions. The
current in P is 3 times the current in Q. The magnetic force on Q has size $F$. Which
statement is correct?
\choices[1]{The wires attract, and the force on P has size $F$.}
  {The wires repel, and the force on P has size $F$.}
  {The wires repel, and the force on P has size $3F$.}
  {The wires repel, and the force on P has size $F/3$.}
\end{practiceq}

\Needspace{11\baselineskip}
\begin{practiceq}{1}{5}\label{pq:induction}
Which statement about a simple AC induction motor is correct?
\choices[1]{Brushes and a split-ring commutator carry the current into the rotor.}
  {The rotor turns at exactly the same speed as the rotating magnetic field.}
  {The rotor turns a little more slowly than the rotating magnetic field.}
  {The rotor must contain a permanent magnet.}
\end{practiceq}

\Needspace{9\baselineskip}
\begin{practiceq}{1}{3}\label{pq:fluxratio}
A flat circular loop of radius $r$ is at right angles to a uniform magnetic field. The
magnetic flux through it is $\Phi$. It is replaced by a loop of radius $2r$ in the same
field. The new loop is turned. Its normal (the line at \ang{90} to the plane of the loop)
now makes an angle of \ang{60} with the field. What is the flux through the new loop?
\choices{$\Phi$}{$2\Phi$}{$2\sqrt{3}\,\Phi$}{$4\Phi$}
\end{practiceq}

\Needspace{15\baselineskip}
\begin{practiceq}{1}{3}\label{pq:ring}
\begin{sidebyside}[0.33]
  \centering
  \begin{tikzpicture}[x=1cm, y=1cm]
    \fill[ground] (-1.45,-0.22) rectangle (1.0,0);
    \draw[hscink, line width=0.6pt] (-1.45,0) -- (1.0,0);
    % iron rod
    \filldraw[fill=hscmuted!30, draw=hscink, line width=0.5pt] (-0.15,0) rectangle (0.15,2.75);
    % solenoid turns
    \foreach \y in {0.2,0.4,...,1.61}
      \draw[hscseriesB!85!black, line width=1pt] (0,\y) ellipse (0.48 and 0.085);
    % aluminium ring
    \draw[hscmuted!75!black, line width=2.4pt] (0,2.2) ellipse (0.34 and 0.075);
    \node[diagramlabel, font=\footnotesize, anchor=west, align=left] at (0.36,2.25)
      {aluminium\\ring};
    \node[diagramlabel, font=\footnotesize, anchor=south] at (0,2.78) {iron rod};
    \node[diagramlabel, font=\footnotesize, anchor=west] at (0.5,0.9) {solenoid};
    % circuit: battery and switch on the left
    \draw[hscink, line width=0.6pt] (-0.48,1.6) -- (-1.15,1.6) -- (-1.15,1.42);
    \fill[hscink] (-1.15,1.42) circle (1.2pt);
    \draw[hscink, line width=0.8pt] (-1.15,1.42) -- (-1.15,1.12);
    \fill[hscink] (-1.15,1.12) circle (1.2pt);
    \node[diagramlabel, font=\scriptsize, anchor=east] at (-1.2,1.27) {switch};
    \draw[hscink, line width=0.6pt] (-1.15,1.12) -- (-1.15,0.86);
    \draw[hscink, line width=1.1pt] (-1.36,0.86) -- (-0.94,0.86);
    \draw[hscink, line width=2pt] (-1.24,0.74) -- (-1.06,0.74);
    \draw[hscink, line width=0.6pt] (-1.15,0.74) -- (-1.15,0.2) -- (-0.48,0.2);
  \end{tikzpicture}
\nextside
A light aluminium ring sits on top of a vertical solenoid. An iron rod passes through the
solenoid and the ring (see the diagram). The switch is closed. Viewed from above, the
current in the solenoid then flows anticlockwise. Just after the switch is closed, which way
does the induced current flow in the ring (viewed from above), and what happens to the
ring?
\end{sidebyside}
\choices[2]{clockwise; the ring is pushed up}
  {clockwise; the ring is pulled down}
  {anticlockwise; the ring is pushed up}
  {anticlockwise; the ring is pulled down}
\end{practiceq}

\Needspace{9\baselineskip}
\begin{practiceq}{1}{1}\label{pq:alpha}
A proton and an alpha particle each start from rest. Each is accelerated through the same
voltage $V$. Each then enters the same uniform magnetic field at \ang{90} to the field. An
alpha particle has 4 times the mass and 2 times the charge of a proton. The radius of the
alpha particle's path is $r_\alpha$. The radius of the proton's path is $r_\mathrm{p}$. What is
the value of $r_\alpha / r_\mathrm{p}$?
\choices{$\dfrac{1}{2}$}{$1$}{$\sqrt{2}$}{$2$}
\end{practiceq}

% =====================================================================
\section*{Section II: Short answer}

\begin{practiceq}{3}{1}\label{pq:gun}
In an electron gun, electrons leave a hot wire with almost no speed. A voltage of
\qty{2.5}{kV} accelerates them across a gap of \qty{1.0}{cm}. Treat the field in the gap as
uniform.
\begin{enumerate}[label=(\alph*)]
  \item Calculate the electric field strength in the gap. \pracmarks{1 mark}
  \item Calculate the kinetic energy that each electron gains, in joules. \pracmarks{1 mark}
  \item Calculate the speed of an electron as it leaves the gap. \pracmarks{1 mark}
\end{enumerate}
\end{practiceq}

\begin{practiceq}{3}{2}\label{pq:busbars}
In a factory switchboard, straight copper bars carry large currents to the machines. (These
bars are called busbars.) Two busbars run side by side for \qty{1.2}{m}. They are parallel
and \qty{4.0}{cm} apart (centre to centre). Each bar carries a current of \qty{500}{A}. The
two currents flow in opposite directions.
\begin{enumerate}[label=(\alph*)]
  \item Calculate the size of the magnetic force on each bar. State whether the bars attract
        or repel each other. \pracmarks{2 marks}
  \item Before 2019, the SI definition of the ampere was based on this kind of force. Two
        long, straight, parallel wires are \qty{1}{m} apart in a vacuum. Each carries a
        current of \qty{1}{A}. State the size of the force on each metre of each wire.
        \pracmarks{1 mark}
\end{enumerate}
\end{practiceq}

\begin{practiceq}{3}{4}\label{pq:neon}
A neon sign needs \qty{7.2}{kV} AC. A transformer steps up a \qty{240}{V} AC supply. Its
primary coil has 300 turns.
\begin{enumerate}[label=(\alph*)]
  \item Calculate the number of turns on the secondary coil. \pracmarks{1 mark}
  \item The current in the secondary coil is \qty{20}{mA}. Assume the transformer is ideal.
        Calculate the current in the primary coil. \pracmarks{1 mark}
  \item In fact, the current in the primary coil is \qty{0.75}{A}. The secondary coil still
        gives \qty{7.2}{kV} and \qty{20}{mA}. Calculate the efficiency of the transformer.
        \pracmarks{1 mark}
\end{enumerate}
\end{practiceq}

\Needspace{13\baselineskip}
\begin{practiceq}{4}{2}\label{pq:rod}
\begin{sidebyside}[0.40]
  \centering
  \begin{tikzpicture}[x=1cm, y=1cm]
    \fill[ground] (-1.75,2.45) rectangle (1.75,2.65);
    \draw[hscink, line width=0.6pt] (-1.75,2.45) -- (1.75,2.45);
    \draw[hscink, line width=0.6pt] (-1.3,2.45) -- (-1.3,0);
    \draw[hscink, line width=0.6pt] (1.3,2.45) -- (1.3,0);
    \draw[guide] (-1.75,-0.6) rectangle (1.75,0.75);
    \foreach \x in {-1.5,-0.9,-0.3,0.3,0.9,1.5}
      \foreach \y in {-0.35,0.45}
        \pic at (\x,\y) {praccross};
    \draw[hscink!80, line width=3pt] (-1.5,0) -- (1.5,0);
    \node[diagramlabel, font=\footnotesize, anchor=north] at (0,-0.62) {copper rod, \qty{0.20}{m}};
    \node[diagramlabel, font=\footnotesize, anchor=west, align=left] at (1.35,1.75)
      {light\\flexible\\wires};
    \node[diagramlabel, font=\footnotesize, text=hscfield, align=center] at (0,1.3)
      {field region:\\$B$ into the page};
  \end{tikzpicture}
\nextside
A copper rod hangs from two light, flexible wires (see the diagram). The rod is horizontal.
It is \qty{0.20}{m} long and has a mass of \qty{8.0}{g}. The rod is in a uniform
\qty{0.050}{T} magnetic field. The field is horizontal, into the page, at \ang{90} to the
rod. A current is passed through the rod. It makes the tension in the two wires zero.
\end{sidebyside}
\begin{enumerate}[label=(\alph*)]
  \item Calculate the current in the rod. \pracmarks{2 marks}
  \item State whether the current flows to the left or to the right. Justify your answer.
        \pracmarks{1 mark}
  \item The rod is now turned in the horizontal plane until it makes an angle of \ang{30}
        with the field. (The rod and the field are both still horizontal.) Calculate the
        current now needed to make the tension in the wires zero. \pracmarks{1 mark}
\end{enumerate}
\end{practiceq}

\begin{practiceq}{4}{3}\label{pq:flip}
A flat circular coil has 75 turns and a diameter of \qty{4.0}{cm}. It sits between the poles
of a large magnet, where the field is uniform and \qty{0.35}{T}. The plane of the coil is at
right angles to the field. In \qty{0.060}{s}, the coil is flipped over (turned through
\ang{180}), so that its plane is again at right angles to the field.
\begin{enumerate}[label=(\alph*)]
  \item Calculate the magnetic flux through one turn of the coil before it is flipped.
        \pracmarks{1 mark}
  \item Calculate the size of the average emf induced in the coil while it is flipped.
        \pracmarks{2 marks}
  \item Apart from turning the coil, state two other ways to change the flux through it.
        \pracmarks{1 mark}
\end{enumerate}
\end{practiceq}

\Needspace{13\baselineskip}
\begin{practiceq}{4}{5}\label{pq:torque}
\begin{sidebyside}[0.42]
  \centering
  \begin{tikzpicture}[x=1cm, y=1cm]
    \draw[pracN] (-2.35,-1.35) rectangle (-2.05,1.35);
    \node[font=\small\bfseries] at (-2.2,1.6) {N};
    \draw[pracS] (2.05,-1.35) rectangle (2.35,1.35);
    \node[font=\small\bfseries] at (2.2,1.6) {S};
    \foreach \y in {-1.1,1.1}
      \draw[field, hscfield!55] (-2.05,\y) -- (2.05,\y);
    \draw[field, hscfield!55] (-2.05,0) -- (-1.35,0);
    \draw[field, hscfield!55] (1.35,0) -- (2.05,0);
    \draw[guide] (-1.3,0) -- (1.3,0);
    \draw[hscink, line width=2pt] (200:1.2) -- (20:1.2);
    \fill[hscink] (0,0) circle (2pt);
    \draw[hscink, line width=0.6pt] (0.9,0) arc (0:20:0.9);
    \node[font=\footnotesize] at (10:1.6) {\ang{20}};
    \node[diagramlabel, font=\footnotesize, anchor=south] at (20:1.25) {coil};
    \node[diagramlabel, font=\footnotesize, anchor=north] at (0,-1.45) {axle (end-on)};
  \end{tikzpicture}
\nextside
The coil of a small DC motor is a rectangle, \qty{3.0}{cm} by \qty{5.0}{cm}. It has 120
turns and carries a current of \qty{0.75}{A}. It sits in a uniform \qty{0.18}{T} field
between two flat magnets. The diagram shows the coil end-on, at one moment as it turns.
\end{sidebyside}
\begin{enumerate}[label=(\alph*)]
  \item Calculate the maximum torque on the coil. \pracmarks{1 mark}
  \item Calculate the torque at the moment shown, when the plane of the coil makes an angle
        of \ang{20} with the magnetic field. \pracmarks{2 marks}
  \item Explain why the motor needs a split-ring commutator to keep turning in one direction.
        \pracmarks{1 mark}
\end{enumerate}
\end{practiceq}

\begin{practiceq}{3}{5}\label{pq:zipline}
A zipline is a ride in which a trolley carries a rider along a long cable. Together the
trolley and the rider have a mass of \qty{90}{kg}. The last part of the line is flat. There,
an aluminium fin (a thin, flat plate) on the trolley passes between rows of strong permanent
magnets. This
magnetic (eddy-current) brake slows the trolley from \qty{15}{m.s^{-1}} to
\qty{3.0}{m.s^{-1}}.
\begin{enumerate}[label=(\alph*)]
  \item Calculate the energy that the brake changes into heat. Ignore friction and air
        resistance. \pracmarks{1 mark}
  \item The brake slows the trolley quickly at first, then more and more gently. Explain why
        the braking force becomes smaller as the trolley slows down. \pracmarks{2 marks}
\end{enumerate}
\end{practiceq}

\begin{practiceq}{4}{4}\label{pq:hydro}
A hydroelectric power station sends out \qty{18}{MW} of electrical power to a town. The
transmission lines have a total resistance of \qty{4.2}{\ohm}.
\begin{enumerate}[label=(\alph*)]
  \item The power is sent at \qty{66}{kV}. Calculate the current in the lines and the power
        lost in them. \pracmarks{2 marks}
  \item The voltage is now stepped up to \qty{330}{kV} before transmission. Calculate the new
        power lost in the lines. \pracmarks{1 mark}
  \item Use your answers to explain why electricity is transmitted at high voltage.
        \pracmarks{1 mark}
\end{enumerate}
\end{practiceq}

\begin{practiceq}{4}{3}\label{pq:fluxgraph}
Coil S has 30 turns. It is wound around the middle of a long solenoid P. The current in P
is changed. This changes the magnetic flux through each turn of S, as shown in the graph.
\begin{figure}
  \centering
  \begin{tikzpicture}
    \begin{axis}[hsc, width=0.80\linewidth, height=0.34\linewidth,
        xmin=0, xmax=100, ymin=-1.6, ymax=1.6, xlabel={$t$ (\si{ms})},
        ylabel={$\Phi$ (\si{mWb})}, xtick={0,10,...,100}, ytick={-1.2,-0.6,0,0.6,1.2},
        yticklabel style={/pgf/number format/.cd, fixed, fixed zerofill, precision=1}]
      \draw[hscink!70, line width=0.6pt] (axis cs:0,0) -- (axis cs:100,0);
      \addplot[hscseriesA, line width=1.3pt, mark=none]
        coordinates {(0,0) (40,1.2) (60,1.2) (80,-1.2) (100,-1.2)};
    \end{axis}
  \end{tikzpicture}
  \caption*{Flux through each turn of coil S.}
\end{figure}
\begin{enumerate}[label=(\alph*)]
  \item Calculate the size of the average emf induced in S between $t = 0$ and
        $t = \qty{40}{ms}$. \pracmarks{1 mark}
  \item A student says: ``At $t = \qty{70}{ms}$ the flux through S is zero, so the emf in S
        is zero.'' Explain why the student is wrong. \pracmarks{1 mark}
  \item On the grid below, sketch the emf induced in S from $t = 0$ to $t = \qty{100}{ms}$.
        Take the emf between $t = 0$ and $t = \qty{40}{ms}$ as negative. Show the values.
        \pracmarks{2 marks}
\end{enumerate}
\begin{figure}
  \centering
  \begin{tikzpicture}
    \begin{axis}[hsc, width=0.80\linewidth, height=0.36\linewidth,
        xmin=0, xmax=100, ymin=-4, ymax=4, xlabel={$t$ (\si{ms})},
        ylabel={$\varepsilon$ (\si{V})}, xtick={0,10,...,100}, ytick={-4,-3,...,4},
        minor y tick num=1, grid=both, minor grid style={hscrule!45, line width=0.2pt}]
      \draw[hscink!70, line width=0.6pt] (axis cs:0,0) -- (axis cs:100,0);
      \addplot[draw=none] coordinates {(0,0)};
    \end{axis}
  \end{tikzpicture}
  \caption*{Grid for part (c).}
\end{figure}
\end{practiceq}

\Needspace{14\baselineskip}
\begin{practiceq}{4}{1}\label{pq:spectro}
\begin{sidebyside}[0.47]
  \centering
  \begin{tikzpicture}[x=0.9cm, y=0.9cm]
    % to scale: radii 2.494 and 2.909 units = 3.118 cm and 3.636 cm (1 unit = 1.25 cm)
    \foreach \x in {-0.2,0.4,...,6.2}
      \foreach \y in {0.4,1.0,...,3.4}
        \pic at (\x,\y) {pracdot};
    \fill[ground] (-0.6,-0.22) rectangle (-0.12,0);
    \fill[ground] (0.12,-0.22) rectangle (6.4,0);
    \draw[hscink, line width=0.8pt] (-0.6,0) -- (-0.12,0);
    \draw[hscink, line width=1.6pt] (0.12,0) -- (6.4,0);
    \draw[velocity] (0,-1.0) -- (0,-0.1);
    \node[diagramlabel, font=\footnotesize, anchor=north] at (0,-1.0) {ions enter};
    \draw[hscseriesA, line width=1pt, midarrow] (0,0) arc (180:0:2.494);
    \draw[hscseriesB, line width=1pt, dashed] (0,0) arc (180:0:2.909);
    \fill[hscink] (4.989,0) circle (2pt) node[below=3pt, font=\small] {P};
    \fill[hscink] (5.818,0) circle (2pt) node[below=3pt, font=\small] {Q};
    \node[diagramlabel, font=\footnotesize, anchor=north] at (3.0,-0.3) {detector plate};
    \node[diagramlabel, font=\footnotesize, text=hscfield, anchor=south] at (3.0,3.6)
      {$B$ out of the page};
  \end{tikzpicture}
\nextside
In a mass spectrometer, lithium ions with a charge of $+e$ all move at
\qty{2.0e5}{m.s^{-1}}. They enter a uniform \qty{0.40}{T} magnetic field at \ang{90} to the
field. Each ion moves in a semicircle and hits a detector plate at P or Q (see the diagram,
drawn to scale). A lithium-6 ion has a mass of \qty{9.99e-27}{kg}. A lithium-7 ion has a
mass of \qty{1.165e-26}{kg}.
\end{sidebyside}
\begin{enumerate}[label=(\alph*)]
  \item Show that the radius of the path of an ion is $r = \dfrac{mv}{qB}$.
        \pracmarks{2 marks}
  \item Calculate the distance PQ. \pracmarks{2 marks}
\end{enumerate}
\end{practiceq}

% =====================================================================
\Needspace{17\baselineskip}
\section*{Section III: Extended response}

\begin{practiceq}{5}{4}\label{pq:evaluate}
A student tests a step-down transformer that is used to charge a power tool. The turns ratio
is $N_\mathrm{p} : N_\mathrm{s} = 12 : 1$. The table shows her results.
\begin{table}
  \caption*{Measurements on the transformer while it charges the tool.}
  \centering
  \begin{tabular}{@{}lcc@{}}
    \toprule
    & \colhead{Voltage (\si{V})} & \colhead{Current (\si{A})}\\
    \midrule
    Primary coil & 240 & 0.275\\
    Secondary coil & 18.6 & 3.20\\
    \bottomrule
  \end{tabular}
\end{table}
The student claims: ``This transformer behaves like an ideal transformer.''

Evaluate the student's claim. In your answer, use the data. Refer to the causes of energy
loss in real transformers and to the ways that designers reduce them.
\end{practiceq}

\Needspace{16\baselineskip}
\begin{practiceq}{6}{5}\label{pq:generator}
The diagram shows a simple generator. A coil turns between two magnets. Each end of the coil
is joined to a slip ring. Brushes connect the slip rings to a lamp. The coil turns at a
steady 25 turns each second. The peak emf is \qty{4.0}{V}.
\begin{figure}
  \centering
  \begin{tikzpicture}[x={(1cm,0cm)}, y={(0cm,1cm)}, z={(0.42cm,0.30cm)}, scale=0.85]
    \def\mz{2.8}
    \filldraw[draw=hscink, line width=0.5pt, fill=hscforce!28] (-2.6,-1,0) -- (-2.6,-1,\mz) -- (-2.6,1,\mz) -- (-2.6,1,0) -- cycle;
    \filldraw[draw=hscink, line width=0.5pt, fill=hscforce!14] (-3.7,1,0) -- (-2.6,1,0) -- (-2.6,1,\mz) -- (-3.7,1,\mz) -- cycle;
    \filldraw[draw=hscink, line width=0.5pt, fill=hscforce!22] (-3.7,-1,0) -- (-2.6,-1,0) -- (-2.6,1,0) -- (-3.7,1,0) -- cycle;
    \node[font=\bfseries] at (-3.15,0,0) {N};
    \draw[field] (-2.6,0.62,1.4) -- (2.6,0.62,1.4);
    \draw[field] (-2.6,-0.4,2.6) -- (2.6,-0.4,2.6);
    \filldraw[draw=hscink, line width=0.5pt, fill=hscvel!14] (2.6,1,0) -- (3.7,1,0) -- (3.7,1,\mz) -- (2.6,1,\mz) -- cycle;
    \filldraw[draw=hscink, line width=0.5pt, fill=hscvel!28] (3.7,-1,0) -- (3.7,-1,\mz) -- (3.7,1,\mz) -- (3.7,1,0) -- cycle;
    \filldraw[draw=hscink, line width=0.5pt, fill=hscvel!22] (2.6,-1,0) -- (3.7,-1,0) -- (3.7,1,0) -- (2.6,1,0) -- cycle;
    \node[font=\bfseries] at (3.15,0,0) {S};
    \draw[hscmuted, line width=0.9pt, dashed] (0,0,-2.4) -- (0,0,3.4);
    \draw[hscink, line width=1.6pt] (0,0,3.4) -- (0,0.75,3.4) -- (0,0.75,3.9);
    \draw[->, hscmuted, line width=0.8pt] ([shift={(0,0,3.4)}]150:0.55) arc (150:30:0.55);
    \node[diagramlabel, anchor=west] at (0.5,1.0,3.6) {handle};
    \draw[pracWire] (1.4,0,0.3) -- (1.4,0,2.3) -- (-1.4,0,2.3) -- (-1.4,0,0.3) -- (-0.12,0,0.3) -- (-0.12,0,-0.4);
    \draw[pracWire] (1.4,0,0.3) -- (0.12,0,0.3) -- (0.12,0,-1.95);
    \draw[hscseriesA, line width=2.6pt] (0,0,-2.0) circle (0.27);
    \draw[hscseriesB, line width=2.6pt] (0,0,-0.45) circle (0.27);
    \filldraw[fill=hscmuted!55, draw=hscink, line width=0.4pt] (-0.07,-0.33,-0.45) rectangle (0.07,-0.62,-0.45);
    \filldraw[fill=hscmuted!55, draw=hscink, line width=0.4pt] (-0.07,-0.33,-2.0) rectangle (0.07,-0.62,-2.0);
    \draw[hscink, line width=0.8pt] (0,-0.62,-0.45) -- (0,-1.0,-0.45) -- (1.6,-1.0,-0.45) -- (1.6,-2.05,-0.45) -- (1.6,-2.05,-1.2) -- (0.39,-2.05,-1.2);
    \draw[hscink, line width=0.8pt] (0,-0.62,-2.0) -- (0,-2.05,-2.0) -- (-0.09,-2.05,-1.2);
    \begin{scope}[shift={(0.15,-2.05,-1.2)}]
      \filldraw[fill=hsctip!25, draw=hscink, line width=0.7pt] (0,0) circle (0.24);
      \draw[hscink, line width=0.6pt] (-0.17,-0.17) -- (0.17,0.17) (-0.17,0.17) -- (0.17,-0.17);
    \end{scope}
    \node[diagramlabel, anchor=north west] at (-0.05,-2.35,-1.2) {lamp};
    \node[diagramlabel, anchor=north east, align=right, inner sep=1pt] at (-1.2,-0.8,0) {slip\\ring};
    \draw[hscmuted, line width=0.4pt] (-0.27,0,-2.0) -- (-1.18,-0.85,0);
    \node[diagramlabel, anchor=west, inner sep=1pt] at (0.62,-0.42,0) {slip ring};
    \draw[hscmuted, line width=0.4pt] (0.27,0,-0.45) -- (0.6,-0.38,0);
    \draw[hscmuted, line width=0.4pt] (0.07,-0.5,-0.45) -- (2.75,-1.25,-0.45) node[anchor=west, diagramlabel, inner sep=1pt] {brushes};
    \draw[hscmuted, line width=0.4pt] (-0.4,0,2.3) -- ++(-0.2,1.05) node[anchor=south, diagramlabel, inner sep=1pt] {coil};
  \end{tikzpicture}
  \caption*{A simple generator with slip rings.}
\end{figure}
\begin{enumerate}[label=(\alph*)]
  \item On the grid below, sketch the output emf for two full turns of the coil. Start when
        the plane of the coil is at right angles to the magnetic field. \pracmarks{2 marks}
  \item Describe how the output would change if:
        (i) the slip rings were replaced by a split-ring commutator;
        (ii) the coil had twice as many turns, and turned at the same speed.
        \pracmarks{2 marks}
  \item The handle is harder to turn when the lamp is connected than when the circuit is
        open. Explain why, using Lenz's law and the law of conservation of energy.
        \pracmarks{2 marks}
\end{enumerate}
\begin{figure}
  \centering
  \begin{tikzpicture}
    \begin{axis}[hsc, width=0.80\linewidth, height=0.36\linewidth,
        xmin=0, xmax=0.08, ymin=-5, ymax=5, xlabel={$t$ (\si{s})},
        ylabel={$\varepsilon$ (\si{V})}, xtick={0,0.01,...,0.081}, ytick={-5,-4,...,5},
        xticklabel style={/pgf/number format/.cd, fixed, fixed zerofill, precision=2},
        grid=both]
      \draw[hscink!70, line width=0.6pt] (axis cs:0,0) -- (axis cs:0.08,0);
      \addplot[draw=none] coordinates {(0,0)};
    \end{axis}
  \end{tikzpicture}
  \caption*{Grid for part (a).}
\end{figure}
\end{practiceq}

\Needspace{13\baselineskip}
\begin{practiceq}{7}{3}\label{pq:coilpull}
A square coil has 25 turns, sides of \qty{10}{cm} and a total resistance of \qty{0.40}{\ohm}.
It is pulled to the right at a steady \qty{2.0}{m.s^{-1}} through a region of uniform
magnetic field. The field is \qty{0.40}{T}, into the page. The region is \qty{30}{cm} wide
(see the diagram). The plane of the coil stays at \ang{90} to the field. At $t = 0$ the
right-hand side of the coil reaches the edge of the field.
\begin{figure}
  \centering
  \begin{tikzpicture}[x=1.15cm, y=1.15cm]
    \draw[hscfield, dashed, line width=0.6pt] (1,-0.5) rectangle (4,1.5);
    \foreach \x in {1.25,1.75,...,3.75}
      \foreach \y in {-0.25,0.25,...,1.25}
        \pic at (\x,\y) {praccross};
    \draw[pracWire] (0,0) rectangle (1,1);
    \node[diagramlabel, font=\footnotesize, anchor=east] at (-0.1,0.78) {coil (25 turns)};
    \draw[velocity] (-0.8,0.3) -- (-0.1,0.3);
    \node[anchor=east, text=hscvel, font=\footnotesize] at (-0.85,0.3)
      {$v = \qty{2.0}{m.s^{-1}}$};
    \draw[hscmuted, <->, line width=0.5pt] (0,-0.25) -- (1,-0.25)
      node[midway, below, font=\footnotesize, text=hscink] {\qty{10}{cm}};
    \draw[hscmuted, <->, line width=0.5pt] (1,1.75) -- (4,1.75)
      node[midway, above, font=\footnotesize, text=hscink] {\qty{30}{cm}};
    \node[diagramlabel, font=\footnotesize, text=hscfield, anchor=west] at (4.1,0.5)
      {$B = \qty{0.40}{T}$, into the page};
  \end{tikzpicture}
  \caption*{The coil at $t = 0$ (seen from the front).}
\end{figure}
\begin{enumerate}[label=(\alph*)]
  \item Calculate the size of the emf induced in the coil while it enters the field.
        \pracmarks{2 marks}
  \item On the grid below, sketch the emf in the coil from $t = 0$ to $t = \qty{0.25}{s}$.
        Take the emf while the coil enters the field as positive. \pracmarks{2 marks}
  \item Calculate the force needed to keep the coil moving at a steady speed while it enters
        the field. Explain, using Lenz's law and energy, why this force is needed.
        \pracmarks{3 marks}
\end{enumerate}
\begin{figure}
  \centering
  \begin{tikzpicture}
    \begin{axis}[hsc, width=0.80\linewidth, height=0.30\linewidth,
        xmin=0, xmax=0.25, ymin=-3, ymax=3, xlabel={$t$ (\si{s})},
        ylabel={$\varepsilon$ (\si{V})}, xtick={0,0.05,...,0.251}, ytick={-3,-2,...,3},
        xticklabel style={/pgf/number format/.cd, fixed, fixed zerofill, precision=2},
        minor x tick num=1, minor y tick num=1, grid=both,
        minor grid style={hscrule!45, line width=0.2pt}]
      \draw[hscink!70, line width=0.6pt] (axis cs:0,0) -- (axis cs:0.25,0);
      \addplot[draw=none] coordinates {(0,0)};
    \end{axis}
  \end{tikzpicture}
  \caption*{Grid for part (b).}
\end{figure}
\end{practiceq}

\Needspace{18\baselineskip}
\begin{practiceq}{6}{2}\label{pq:rods}
A student investigates the force between two parallel conductors. Two straight aluminium rods
are each \qty{0.60}{m} long. The lower rod lies on insulating blocks on an electronic
balance. The upper rod is clamped directly above it. The rods are joined in series, so the
same current flows in both, in opposite directions. The current is kept at \qty{8.0}{A}.
The student changes $r$, the distance between the centres of the rods. To do this, she
uses thin insulating spacers (small blocks of known thickness). With no current, the
balance is set to zero. The balance then shows the extra reading $\Delta m$ caused by the
magnetic force.
\begin{figure}
  \centering
  \begin{tikzpicture}[x=1cm, y=1cm]
    % balance
    \draw[object] (0.9,-1.15) rectangle (5.1,-0.35);
    \draw[hscink, line width=0.6pt, fill=white] (2.3,-0.95) rectangle (3.7,-0.55);
    \node[font=\scriptsize] at (3.0,-0.75) {$\Delta m$ (g)};
    \draw[hscink, line width=0.8pt, fill=hscmuted!25] (0.3,-0.35) rectangle (5.7,-0.22);
    \fill[hscmuted!60] (1.0,-0.22) rectangle (1.3,0.07);
    \fill[hscmuted!60] (4.7,-0.22) rectangle (5.0,0.07);
    \node[diagramlabel, font=\footnotesize, anchor=north] at (3.0,-1.2) {electronic balance};
    % rods (separation exaggerated); lower rod centre y = 0.12, upper rod centre y = 1.12
    \draw[hscink!75, line width=3pt] (0,0.12) -- (6,0.12);
    \draw[hscink!75, line width=3pt] (0,1.12) -- (6,1.12);
    \draw[->, hscink, line width=1pt] (2.3,1.42) -- (3.7,1.42) node[right, font=\small] {$I$};
    \draw[<-, hscink, line width=1pt] (2.3,0.4) -- (3.7,0.4) node[right, font=\small] {$I$};
    % clamps for the upper rod
    \draw[hscink, line width=0.8pt] (0.6,1.12) -- (0.6,2.0);
    \draw[hscink, line width=0.8pt] (5.4,1.12) -- (5.4,2.0);
    \fill[ground] (0.2,2.0) rectangle (5.8,2.18);
    \draw[hscink, line width=0.6pt] (0.2,2.0) -- (5.8,2.0);
    % series link at the right ends
    \draw[hscink, line width=0.7pt] (6,1.12) .. controls (6.7,1.12) and (6.7,0.12) .. (6,0.12);
    % supply at the left
    \draw[hscink, line width=0.7pt] (0,1.12) -- (-0.8,1.12) -- (-0.8,1.5);
    \draw[hscink, line width=0.7pt] (0,0.12) -- (-1.6,0.12) -- (-1.6,1.5);
    \draw[object] (-2.3,1.5) rectangle (-0.1,2.3);
    \node[font=\scriptsize, align=center] at (-1.2,1.9) {DC supply\\and ammeter};
    % separation between the centres
    \draw[hscmuted, <->, line width=0.5pt] (7.0,0.12) -- (7.0,1.12);
    \node[font=\footnotesize, anchor=west] at (7.05,0.62) {$r$};
  \end{tikzpicture}
  \caption*{The set-up (not to scale).}
\end{figure}
\begin{table}
  \caption*{Results ($I = \qty{8.0}{A}$ in each rod).}
  \centering
  \begin{tabular}{@{}lccccc@{}}
    \toprule
    Distance $r$ (\si{mm}) & 3.0 & 4.0 & 5.0 & 6.0 & 8.0\\
    \midrule
    $\dfrac{1}{r}$ (\si{m^{-1}}) & \pracblank & \pracblank & \pracblank & \pracblank
      & \pracblank\\
    \midrule
    Balance reading $\Delta m$ (\si{g}) & 0.256 & 0.193 & 0.153 & 0.129 & 0.096\\
    \midrule
    Force $F$ ($\times 10^{-3}$ \si{N}) & \pracblank & \pracblank & \pracblank & \pracblank
      & \pracblank\\
    \bottomrule
  \end{tabular}
\end{table}
\begin{enumerate}[label=(\alph*)]
  \item Complete the table. \pracmarks{1 mark}
  \item On the grid for part (b), plot $F$ against $\dfrac{1}{r}$. Draw a line of best fit.
        \pracmarks{2 marks}
  \item Use the gradient of your line to calculate a value for $\mu_0$. \pracmarks{2 marks}
  \item The formula for the force between two wires is for very long, straight wires.
        Explain why the rods must be much longer than the distance between them for the
        formula to apply. \pracmarks{1 mark}
\end{enumerate}
\begin{figure}
  \centering
  \begin{tikzpicture}
    \begin{axis}[hsc, width=0.84\linewidth, height=0.56\linewidth,
        xmin=0, xmax=350, ymin=0, ymax=3.0,
        xlabel={$\dfrac{1}{r}$ (\si{m^{-1}})}, ylabel={$F$ ($\times 10^{-3}$ \si{N})},
        xtick={0,50,...,350}, ytick={0,0.5,...,3.0},
        minor x tick num=1, minor y tick num=4,
        grid=both, minor grid style={hscrule!45, line width=0.2pt},
        yticklabel style={/pgf/number format/.cd, fixed, fixed zerofill, precision=1}]
      \addplot[draw=none] coordinates {(0,0)};
    \end{axis}
  \end{tikzpicture}
  \caption*{Grid for part (b).}
\end{figure}
\end{practiceq}

\Needspace{15\baselineskip}
\begin{practiceq}{6}{5}\label{pq:motorgraph}
A small DC motor is connected to a \qty{6.0}{V} DC supply. It is switched on at $t = 0$. The
graph shows the current in the motor. At $t = \qty{2.0}{s}$ a load is attached to the
motor's axle. (A load is something that the motor must turn, so the motor must now do work.)
\begin{figure}
  \centering
  \begin{tikzpicture}
    \begin{axis}[hsc, width=0.84\linewidth, height=0.46\linewidth,
        xmin=0, xmax=3.0, ymin=0, ymax=4.4, xlabel={$t$ (\si{s})},
        ylabel={$I$ (\si{A})}, xtick={0,0.5,...,3.0}, ytick={0,1,...,4},
        minor y tick num=4, grid=both, minor grid style={hscrule!45, line width=0.2pt},
        xticklabel style={/pgf/number format/.cd, fixed, fixed zerofill, precision=1},
        yticklabel style={/pgf/number format/.cd, fixed, fixed zerofill, precision=1}]
      \addplot[hscseriesA, line width=1.3pt, mark=none, domain=0:2, samples=161]
        {0.40 + 3.60*exp(-x/0.25)};
      \addplot[hscseriesA, line width=1.3pt, mark=none, domain=2:3, samples=81]
        {1.00 - 0.598792*exp(-(x-2)/0.15)};
      \draw[->, hscmuted, line width=0.7pt] (axis cs:1.75,2.3) -- (axis cs:1.98,0.55);
      \node[anchor=south, font=\footnotesize, text=hscmuted, fill=white, inner sep=1.5pt]
        at (axis cs:1.75,2.3) {load attached};
    \end{axis}
  \end{tikzpicture}
  \caption*{Current in the motor from the moment it is switched on.}
\end{figure}
\begin{enumerate}[label=(\alph*)]
  \item Use the graph to calculate the resistance of the motor's coil. \pracmarks{1 mark}
  \item Calculate the back emf when the motor runs at a steady speed with no load (from
        $t = \qty{1.5}{s}$ to $t = \qty{2.0}{s}$). \pracmarks{1 mark}
  \item Explain the shape of the graph from $t = 0$ to $t = \qty{2.0}{s}$.
        \pracmarks{3 marks}
  \item Explain why the current increases after the load is attached. \pracmarks{1 mark}
\end{enumerate}
\end{practiceq}

\Needspace{13\baselineskip}
\begin{practiceq}{8}{1}\label{pq:crossed}
\begin{sidebyside}[0.46]
  \centering
  \begin{tikzpicture}[x=0.8cm, y=0.8cm]
    \fill[hscforce!20, draw=hscink, line width=0.6pt] (0,1.0) rectangle (5.0,1.15);
    \fill[hscvel!20, draw=hscink, line width=0.6pt] (0,-1.15) rectangle (5.0,-1.0);
    \foreach \x in {0.5,1.5,...,4.5} {
      \node[font=\small\bfseries, text=hscforce] at (\x,1.38) {$+$};
      \node[font=\small\bfseries, text=hscvel] at (\x,-1.38) {$-$};
    }
    \draw[guide] (-1.3,0) -- (5.0,0);
    \fill[object] (-1.2,0) circle (4pt);
    \node[font=\scriptsize\bfseries] at (-1.2,0) {$-$};
    \draw[velocity] (-0.95,0) -- (-0.1,0) node[above, text=hscvel, font=\small] {$v$};
    \draw[hscmuted, <->, line width=0.5pt] (0,-1.85) -- (5.0,-1.85)
      node[midway, below, font=\footnotesize, text=hscink] {$L = \qty{5.0}{cm}$};
    \draw[hscmuted, <->, line width=0.5pt] (5.4,-1.0) -- (5.4,1.0)
      node[midway, right, font=\footnotesize, text=hscink, align=left] {$d =$\\\qty{2.0}{cm}};
  \end{tikzpicture}
\nextside
Two horizontal metal plates are \qty{5.0}{cm} long and \qty{2.0}{cm} apart. The voltage
between them is \qty{120}{V}. The top plate is positive. An electron enters the space
between the plates, halfway between them. It moves parallel to the plates at
\qty{3.0e7}{m.s^{-1}}.
\end{sidebyside}
\begin{enumerate}[label=(\alph*)]
  \item Calculate the vertical deflection of the electron when it leaves the plates. State
        its direction. \pracmarks{3 marks}
  \item A uniform magnetic field is now added in the space between the plates. The electron
        now goes straight through, without any deflection. Calculate the size of the
        magnetic field, and state its direction. \pracmarks{2 marks}
  \item The voltage is now switched off, but the magnetic field stays. Calculate the radius
        of the electron's path. Compare the motion of the electron in this magnetic field
        with its motion in the electric field alone (part (a)). \pracmarks{3 marks}
\end{enumerate}
\end{practiceq}
